% questions.tex —— 1981 年·全国卷（文科） · 题目内容片段（无答案）
% 由 tools/generate_projects.py 依来源总账自 1981/1981_数学真题_文_试题.md 生成，勿手改（改源 MD 后重跑）


\gkpart{文科数学}

\begin{gktip}
\textbf{来源说明}：本题面依据固定原卷 PDF《1981 文》第 1 页（印刷页码 18）逐题转录。源卷为单页三栏题面，未附答案页；配套答案与解析由本库独立推导，\textbf{非官方标准答案}。原卷三栏排布，本文按题号重排为第 1—9 题。原卷全卷无插图。第 7 题附所给对数表，按原卷三列照录。第 2 题含根号与分数，\texttt{pdftotext} 纯文本会把根号 $\sqrt3$、$\sqrt b$ 与分数线错位；本文以页面图像为准转录。
\end{gktip}

\par\medskip

\begin{question}
\gkti{1}设 $A$ 表示有理数的集合，$B$ 表示无理数的集合，即设 $A=\{\text{\text{有理数}}\}$，$B=\{\text{\text{无理数}}\}$，试写出：

（1） $A\cup B$；

（2） $A\cap B$。

\end{question}
\begin{question}
\gkti{2}化简：

\nopagebreak[4]
\[
\left[-\frac{a^7b^2}{\sqrt{3}(a+b)^2}\right]^2\times\left[\frac{a^2-b^2}{a^2\sqrt{b}}\right]^4\div\left[\frac{a^2(b-a)}{2}\right]^3.
\]

\end{question}
\begin{question}
\gkti{3}在 $A$、$B$、$C$、$D$ 四位候选人中，

（1） 如果选举正、副班长各一人，共有几种选法？写出所有可能的选举结果；

（2） 如果选举班委三人，共有几种选法？写出所有可能的选举结果。

\end{question}
\begin{question}
\gkti{4}求函数 $f(x)=\sin x+\cos x$ 在区间 $(-\pi,\pi)$ 上的最大值。

\end{question}
\begin{question}
\gkti{5}写出正弦定理，并对钝角三角形的情况加以证明。

\end{question}
\begin{question}
\gkti{6}已知正方形 $ABCD$ 的相对顶点 $A(0,-1)$ 和 $C(2,5)$，求顶点 $B$ 和 $D$ 的坐标。

\end{question}
\begin{question}
\gkti{7}设 1980 年底我国人口以 10 亿计算。

（1） 如果我国人口每年比上年平均递增 $2\%$，那么到 2000 年底将达到多少？

（2） 要使 2000 年底我国人口不超过 12 亿，那么每年比上年平均递增率最高是多少？

下列对数值可供选用：

\begin{center}
\begin{tabular}{lll}
\toprule
 &  &  \\
\midrule
$\lg 1.0087=0.00377$ & $\lg 1.0092=0.00396$ & $\lg 1.0096=0.00417$ \\
$\lg 1.0200=0.00860$ & $\lg 1.2000=0.07918$ & $\lg 1.3098=0.11720$ \\
$\lg 1.4568=0.16340$ & $\lg 1.4859=0.17200$ & $\lg 1.5157=0.18060$ \\
\bottomrule
\end{tabular}
\end{center}

\end{question}
\begin{question}
\gkti{8}$ABCD-A_1B_1C_1D_1$ 为一正四棱柱，过 $A$、$C$、$B_1$ 三点作一截面，求证：截面 $ACB_1\perp$ 对角面 $DBB_1D_1$。

\end{question}
\begin{question}
\gkti{9}

（1） 设抛物线 $y^2=4x$ 截直线 $y=2x+k$ 所得的弦长为 $3\sqrt5$，求 $k$ 的值。

（2） 以本题（1）得到的弦为底边，以 $x$ 轴上的点 $P$ 为顶点做成三角形。当这三角形的面积为 9 时，求 $P$ 的坐标。

\end{question}